% Bagian: first-order-logic
% Bab: natural-deduction
% Bagian: proving-things-quant

\documentclass[../../../include/open-logic-section]{subfiles}

\begin{document}

\olfileid[id]{fol}{ntd}{prq}

\olsection{\usetoken{P}{derivation} dengan Kuantor}

\begin{ex}
Ketika menangani kuantor, kita harus memastikan bahwa syarat variabel
eigen tidak dilanggar, dan kadang-kadang hal ini mengharuskan kita
mengubah urutan pelaksanaan inferensi tertentu. Secara umum, sebaiknya
kita mencoba menangani terlebih dahulu aturan-aturan yang tunduk pada
syarat variabel eigen (aturan-aturan tersebut akan berada lebih bawah
dalam pembuktian akhir).

Mari kita lihat cara memberikan !!a{derivation} atas !!{formula}
$\lexists[x][\lnot !A(x)] \lif \lnot \lforall[x][!A(x)]$.
Seperti biasa, kita mulai dengan menulis
\begin{prooftree}
\AxiomC{}
\UnaryInfC{$\lexists[x][\lnot !A(x)]\lif \lnot \lforall[x][!A(x)]$}
\end{prooftree}
Kita mulai dengan menuliskan apa yang diperlukan untuk membenarkan
langkah terakhir itu menggunakan aturan \Intro{\lif}.
\begin{prooftree}
\AxiomC{$\Discharge{\lexists[x][\lnot !A(x)]}{1}$}
\DeduceC{$\lnot \lforall[x][!A(x)]$}
\DischargeRule{\Intro{\lif}}{1}
\UnaryInfC{$\lexists[x][\lnot !A(x)]\lif \lnot \lforall[x][!A(x)]$}
\end{prooftree}
Karena tidak ada aturan yang jelas untuk diterapkan pada $\lnot
\lforall[x][!A(x)]$, kita akan menyiapkan !!{derivation} agar dapat
menggunakan aturan \Elim{\lexists}. Di sini kita harus memperhatikan
syarat variabel eigen dan memilih simbol konstanta yang tidak muncul
dalam premis eksistensial $\lexists[x][\lnot !A(x)]$, kesimpulan, ataupun
asumsi yang belum dilepaskan dan menjadi sandarannya. (Namun, karena tidak
ada !!{constant}s yang muncul, pilihan apa pun dapat digunakan.)
\begin{prooftree}
\AxiomC{$\Discharge{\lexists[x][\lnot !A(x)]}{1}$}
\AxiomC{$\Discharge{\lnot !A(a)}{2}$}
\DeduceC{$\lnot \lforall[x][!A(x)]$}
\DischargeRule{\Elim{\lexists}}{2}
\BinaryInfC{$\lnot \lforall[x][!A(x)]$}
\DischargeRule{\Intro{\lif}}{1}
\UnaryInfC{$\lexists[x][\lnot !A(x)] \lif \lnot \lforall[x][!A(x)]$}
\end{prooftree}
Untuk menurunkan $\lnot \lforall[x][!A(x)]$, kita akan mencoba
menggunakan aturan \Intro{\lnot}: aturan ini mengharuskan kita
menurunkan suatu kontradiksi, mungkin dengan menggunakan
$\lforall[x][!A(x)]$ sebagai asumsi tambahan. Tentu saja, kontradiksi
ini boleh melibatkan asumsi $\lnot !A(a)$ yang akan dilepaskan oleh
inferensi \Elim{\lexists}. Kita dapat menyiapkannya sebagai berikut:
\begin{prooftree}
\AxiomC{$\Discharge{\lexists[x][\lnot !A(x)]}{1}$}
\AxiomC{$\Discharge{\lnot !A(a)}{2}, \Discharge{\lforall[x][!A(x)]}{3}$}
\DeduceC{$\lfalse$}
\DischargeRule{\Intro{\lnot}}{3}
\UnaryInfC{$\lnot \lforall[x][!A(x)]$}
\DischargeRule{\Elim{\lexists}}{2}
\BinaryInfC{$\lnot \lforall[x][!A(x)]$}
\DischargeRule{\Intro{\lif}}{1}
\UnaryInfC{$\lexists[x][\lnot !A(x)]\lif \lnot \lforall[x][!A(x)]$}
\end{prooftree}
Tampaknya kita hampir memperoleh kontradiksi. Aturan yang paling mudah
diterapkan adalah \Elim{\lforall}, yang tidak memiliki syarat variabel
eigen. Karena kita dapat menggunakan sebarang suku tertutup untuk
menggantikan $x$ yang terikat kuantor universal, paling masuk akal jika
kita terus menggunakan~$a$ agar dapat mencapai kontradiksi.
\begin{prooftree}
\AxiomC{$\Discharge{\lexists[x][\lnot !A(x)]}{1}$}
\AxiomC{$\Discharge{\lnot !A(a)}{2}$}
\AxiomC{$\Discharge{\lforall[x][!A(x)]}{3}$}
\RightLabel{\Elim{\lforall}}
\UnaryInfC{$!A(a)$}
\RightLabel{\Elim{\lnot}}
\BinaryInfC{$\lfalse$}
\DischargeRule{\Intro{\lnot}}{3}
\UnaryInfC{$\lnot \lforall[x][!A(x)]$}
\DischargeRule{\Elim{\lexists}}{2}
\BinaryInfC{$\lnot \lforall[x][!A(x)]$}
\DischargeRule{\Intro{\lif}}{1}
\UnaryInfC{$\lexists[x][\lnot !A(x)]\lif \lnot \lforall[x][!A(x)]$}
\end{prooftree}

Terutama ketika menangani kuantor, pada tahap ini penting untuk
memeriksa kembali bahwa syarat variabel eigen tidak dilanggar. Karena
satu-satunya aturan yang kita terapkan dan tunduk pada syarat variabel
eigen adalah \Elim{\lexists}, sedangkan variabel eigen~$a$ tidak muncul
dalam premis eksistensial, kesimpulan inferensi, atau asumsi yang belum
dilepaskan, ini merupakan !!a{derivation} yang benar.
\end{ex}

\begin{ex}
Kadang-kadang kita dapat menurunkan !!a{formula} dari !!{formula}s lain.
Dalam kasus seperti ini, mungkin terdapat asumsi yang belum dilepaskan.
Penting untuk tetap mencatat asumsi-asumsi kita selain sasaran akhirnya.

Mari kita lihat cara memberikan !!a{derivation} atas !!{formula}
$\lexists[x][!C(x,b)]$ dari asumsi-asumsi $\lexists[x][(!A(x)
\land !B(x))]$ dan $\lforall[x][(!B(x) \lif !C(x,b))]$.
Seperti biasa, kita mulai dengan menuliskan kesimpulan di bagian bawah.
\begin{prooftree}
\AxiomC{}
\UnaryInfC{$\lexists[x][!C(x,b)]$}
\end{prooftree}

Kita memiliki dua asumsi untuk digunakan. Untuk menggunakan yang
pertama, yakni mencoba menemukan !!a{derivation} atas
$\lexists[x][!C(x, b)]$ dari $\lexists[x][(!A(x) \land !B(x))]$, kita
akan memakai aturan \Elim{\lexists}. Karena aturan ini memiliki syarat
variabel eigen, kita akan menerapkannya terlebih dahulu. Kita memperoleh:
\begin{prooftree}
\AxiomC{$\lexists[x][(!A(x) \land !B(x))]$}
\AxiomC{$\Discharge{!A(a) \land !B(a)}{1}$}
\DeduceC{$\lexists[x][!C(x,b)]$}
\DischargeRule{\Elim{\lexists}}{1}
\BinaryInfC{$\lexists[x][!C(x,b)]$}
\end{prooftree}
Kedua asumsi yang kita gunakan sama-sama melibatkan~$!B$. Pada tahap
ini, mungkin berguna untuk menerapkan \Elim{\land} guna memisahkan
$!B(a)$.
\begin{prooftree}
\AxiomC{$\lexists[x][(!A(x) \land !B(x)])$}
\AxiomC{$\Discharge{!A(a)
\land !B(a)}{1}$}
\RightLabel{\Elim{\land}}
\UnaryInfC{$!B(a)$}
\DeduceC{$\lexists[x][!C(x,b)]$}
\DischargeRule{\Elim{\lexists}}{1}
\BinaryInfC{$\lexists[x][!C(x,b)]$}
\end{prooftree}

Asumsi kedua yang dapat kita gunakan ialah~$\lforall[x][(!B(x) \lif
!C(x,b))]$. Karena tidak ada syarat variabel eigen, kita dapat
menginstansiasi $x$ dengan !!{constant}~$a$ menggunakan
\Elim{\lforall} untuk memperoleh $!B(a) \lif !C(a, b)$. Sekarang kita
memiliki $!B(a) \lif !C(a,b)$ dan $!B(a)$. Langkah selanjutnya adalah
penerapan langsung aturan \Elim{\lif}.
\begin{prooftree}
\AxiomC{$\lexists[x][(!A(x) \land !B(x))]$}
\AxiomC{$\lforall[x][(!B(x) \lif !C(x,b))]$}
\RightLabel{\Elim{\lforall}}
\UnaryInfC{$!B(a) \lif !C(a,b)$}
\AxiomC{$\Discharge{!A(a)
\land !B(a)}{1}$}
\RightLabel{\Elim{\land}}
\UnaryInfC{$!B(a)$}
\RightLabel{\Elim{\lif}}
\BinaryInfC{$!C(a,b)$}
\DeduceC{$\lexists[x][!C(x,b)]$}
\DischargeRule{\Elim{\lexists}}{1}
\insertBetweenHyps{\hspace{-5em}}
\BinaryInfC{$\lexists[x][!C(x,b)]$}
\end{prooftree}
Kita sudah sangat dekat!{} Dengan satu penerapan \Intro{\lexists},
kita mencapai sasaran.
\begin{prooftree}
\AxiomC{$\lexists[x][(!A(x) \land !B(x))]$}
\AxiomC{$\lforall[x][(!B(x) \lif !C(x,b))]$}
\RightLabel{\Elim{\lforall}}
\UnaryInfC{$!B(a) \lif !C(a,b)$}
\AxiomC{$\Discharge{!A(a)
\land !B(a)}{1}$}
\RightLabel{\Elim{\land}}
\UnaryInfC{$!B(a)$}
\RightLabel{\Elim{\lif}}
\BinaryInfC{$!C(a,b)$}
\RightLabel{\Intro{\lexists}}
\UnaryInfC{$\lexists[x][!C(x,b)]$}
\DischargeRule{\Elim{\lexists}}{1}
\insertBetweenHyps{\hspace{-5em}}
\BinaryInfC{$\lexists[x][!C(x,b)]$}
\end{prooftree}
Karena pada setiap langkah kita memastikan bahwa syarat variabel eigen
tidak dilanggar, kita dapat yakin bahwa ini merupakan !!a{derivation}
yang benar.
\end{ex}

\begin{ex}
Berikan !!a{derivation} atas !!{formula} $\lnot\lforall[x][!A(x)]$ dari
asumsi-asumsi $\lforall[x][!A(x)] \lif \lexists[y][!B(y)]$ dan
$\lnot\lexists[y][!B(y)]$.
Seperti biasa, kita menuliskan !!{formula} sasaran di bagian bawah.
\begin{prooftree}
\AxiomC{}
\UnaryInfC{$\lnot\lforall[x][!A(x)]$}
\end{prooftree}
Baris terakhir !!{derivation} merupakan negasi, sehingga mari kita coba
menggunakan \Intro{\lnot}. Ini mengharuskan kita menentukan cara
!!{derive} suatu kontradiksi.
\begin{prooftree}
\AxiomC{$\Discharge{\lforall[x][!A(x)]}{1}$}
\DeduceC{$\lfalse$}
\DischargeRule{\Intro{\lnot}}{1}
\UnaryInfC{$\lnot\lforall[x][!A(x)]$}
\end{prooftree}
Sejauh ini semuanya berjalan baik. Kita dapat menggunakan
\Elim{\lforall}, tetapi belum jelas apakah itu akan membantu mencapai
sasaran. Sebagai gantinya, mari kita gunakan salah satu asumsi.
$\lforall[x][!A(x)] \lif \lexists[y][!B(y)]$ bersama dengan
$\lforall[x][!A(x)]$ memungkinkan kita menggunakan aturan \Elim{\lif}.
\begin{prooftree}
\AxiomC{$\lforall[x][!A(x)] \lif \lexists[y][!B(y)]$}
\AxiomC{$\Discharge{\lforall[x][!A(x)]}{1}$}
\RightLabel{\Elim{\lif}}
\BinaryInfC{$\lexists[y][!B(y)]$}
\DeduceC{$\lfalse$}
\DischargeRule{\Intro{\lnot}}{1}
\UnaryInfC{$\lnot\lforall[x][!A(x)]$}
\end{prooftree}
Kini tinggal satu asumsi lagi untuk digunakan, dan tampaknya asumsi ini
akan membantu kita mencapai kontradiksi dengan menggunakan
\Elim{\lnot}.
\begin{prooftree}
\AxiomC{$\lnot\lexists[y][!B(y)]$}
\AxiomC{$\lforall[x][!A(x)] \lif \lexists[y][!B(y)]$}
\AxiomC{$\Discharge{\lforall[x][!A(x)]}{1}$}
\RightLabel{\Elim{\lif}}
\BinaryInfC{$\lexists[y][!B(y)]$}
\RightLabel{\Elim{\lnot}}
\BinaryInfC{$\lfalse$}
\DischargeRule{\Intro{\lnot}}{1}
\UnaryInfC{$\lnot\lforall[x][!A(x)]$}
\end{prooftree}
\end{ex}

\begin{prob}
Berikan !!{derivation}s yang menunjukkan hal-hal berikut:
\begin{enumerate}
\item $\Proves (\lforall[x][!A(x)] \land \lforall[y][!B(y)]) \lif
\lforall[z][(!A(z) \land !B(z))]$.
\item $\Proves (\lexists[x][!A(x)] \lor \lexists[y][!B(y)]) \lif
\lexists[z][(!A(z) \lor !B(z))]$.
\item $\lforall[x][(!A(x) \lif !B)] \Proves \lexists[y][!A(y)] \lif !B$.
\item $\lforall[x][\lnot !A(x)] \Proves \lnot\lexists[x][!A(x)]$.
\item $\Proves \lnot\lexists[x][!A(x)] \lif \lforall[x][\lnot !A(x)]$.
\item $\Proves \lnot\lexists[x][\lforall[y][((!A(x,y) \lif \lnot
!A(y,y)) \land (\lnot !A(y,y) \lif !A(x,y)))]]$.
\end{enumerate}
\end{prob}

\begin{prob}
Berikan !!{derivation}s yang menunjukkan hal-hal berikut:
\begin{enumerate}
\item $\Proves \lnot\lforall[x][!A(x)] \lif \lexists[x][\lnot!A(x)]$.
\item $(\lforall[x][!A(x)] \lif !B) \Proves \lexists[y][(!A(y) \lif !B)]$.
\item $\Proves \lexists[x][(!A(x) \lif \lforall[y][!A(y)])]$.
\end{enumerate}
(Semua ini memerlukan aturan~$\FalseCl$.)
\end{prob}

\end{document}
